\scanpage{30}
On dit que \deleted{$(A,\Phi)$} \deleted{\unsure{$\varphi$|$\Phi$}} $\varphi$ est \deleted{\unsure{$\varphi$}} \uline{\unsure{riemannien|riemannienne}} si $A$,
muni de sa trace et de l'involution $u\leadsto u^\varphi$, est
une Algèbre riemannienne, i.e. satisfait
\[
\mathrm{Tr}\,u^\varphi u\geq0\qquad \text{tt }u
\]
\deleted{ce qui ne dépend pas du choix de $\varphi$}

[\uline{NB} utilisant les identifications $A\simeq\Phi$, $A^\circ\simeq\Phi^*$
définies par $\varphi$ et $\varphi^{-1}$, la forme $\mathrm{Tr}\,v^\varphi u$
sur $A^\circ\times A$ s'identifie à la forme canonique
sur $\Phi\times\Phi^*$, elle est donc \uline{non dégénérée}]

Tout $\varphi$ riemannienne définit comme plus
haut un cône convexe fermé $\Phi^+(\varphi)=\Phi^+$,
savoir le plus petit cône convexe fermé
stable par $u\otimes u$ ($u\in A$) contenant $\varphi$, et qui
correspond par $\Theta_\varphi:A\simeq\unsure{\Phi|P}$ au cône des
``éléments positifs'' de l'algèbre \deleted{hilb.} riemannienne
$A$, engendré par les éléments de la forme
$u^\varphi u$. De $=$, $A^\circ\longrightarrow\Phi^*$ transforme ce dernier
cône en le cône $\Phi^{*+}$ défini en termes
de $\varphi^{-1}$. La théorie des Alg. hilbertiennes
nous dit que ces cônes $\Phi^{\deleted{\unsure{*}}+}$ et $\unsure{\Phi|\varphi|P}^{*+}$ sont
polaires l'un de l'autre, \unsure{ts.|tq.} que leurs
% archive sheet number: 29
